\documentclass[11pt]{amsart}

\usepackage[T1]{fontenc}
\usepackage{lmodern}
\usepackage{microtype}
\usepackage{mathtools,amssymb,amsthm}
\usepackage[hidelinks]{hyperref}

\hypersetup{
  pdftitle={Voltage assignments on a 44-vertex reconstruction of H12: girth 12 for m = 7, 8, 9, and none for 3 <= m <= 6}
}

\newtheorem{theorem}{Theorem}
\newtheorem{lemma}[theorem]{Lemma}
\newtheorem{proposition}[theorem]{Proposition}
\newtheorem{corollary}[theorem]{Corollary}
\theoremstyle{remark}
\newtheorem*{remark}{Remark}

\title[A $44$-vertex reconstruction of $H_{12}$, for $m\le 9$]
{Voltage assignments on a $44$-vertex reconstruction of $H_{12}$:
 girth~$12$ for $m=7,8,9$, and none for $3\le m\le 6$}

\date{}

\subjclass[2020]{05C35, 05C25, 05C38}
\keywords{bi-regular cage, semicubic graph, girth, voltage graph, lift,
 generalized hexagon, exhaustive search}

\begin{document}

\begin{abstract}
Aguilar, Araujo-Pardo and Berman exhibit a voltage graph $H_{12}$ on $44$ vertices whose
$\mathbb{Z}_m$ lifts are $(\{3,m\};12)$-graphs of order $41m+3$ for every $m\ge 10$, and ask,
in the third sub-bullet under ``Related to semicubic graphs'' of their Section~6 (p.~29), for
voltage assignments on that same $H_{12}$ producing girth~$12$ for $m\in\{3,4,\dots,9\}$, or a
proof that none exists. This note is about one explicit $44$-vertex voltage graph, printed in
full in Section~\ref{sec:base}, which is our reading of the source's Figure~16 and its
Tables~3--4; we do \emph{not} verify that it is the source's $H_{12}$, and we therefore do not
claim to settle the cited question --- whether we do turns entirely on that identification,
which is stated and bounded in Section~\ref{sec:corr}. For the graph printed below: girth-$12$
assignments exist for $m=7,8,9$ --- exhibited here in full, at orders $290$, $331$ and $372$ ---
and none exists for $m=3,4,5,6$. The two smaller witnesses are checked as graphs in their own
right, so they give
\[
 n(\{3,7\};12)\le 290\ \ (\text{was }334),
 \qquad
 n(\{3,8\};12)\le 331\ \ (\text{was }374),
\]
independently of the identification; both lie strictly above the published lower bounds $272$
and $308$. The $m=9$ witness carries \emph{no} record claim: Goedgebeur, Jooken and
Van~den~Eede publish $n(\{3,9\};12)\le 360<372$. The second question of the same sub-bullet,
which asks for a \emph{different} graph $H'_{12}$, is untouched and remains open at every $m$.
\end{abstract}

\maketitle

\section{The question}\label{sec:q}

Following Aguilar, Araujo-Pardo and Berman \cite{AAB}, a \emph{$(\{3,m\};g)$-graph} is a graph
whose every vertex has degree $3$ or $m$ and whose girth is $g$; $n(\{3,m\};g)$ denotes the
least order of such a graph. In a voltage graph over $\mathbb{Z}_m$ a \emph{pinned} vertex
$v^{*}$ is a base vertex of degree~$1$ whose lift is a \emph{single} vertex joined to all $m$
copies of its unique neighbour, so that $\deg v^{*}=m$; a base whose other vertices have
degree~$3$ and which carries $s$ pinned vertices therefore lifts to a graph with exactly $s$
vertices of degree $m$ and all others of degree~$3$ \cite[Remark~3.1]{AAB}. Unlabelled edges
carry voltage~$0$.

The source's Theorem~5.6 (p.~27) reads, verbatim:

\begin{quote}
``The graph $(H_{12};m)$ given by the voltage graph $H_{12}$ shown in Figure 16 with voltages
given in Table 4 has girth 12 for $m \ge 10$. Moreover, it is a $(\{3,m\};12)$-graph with three
vertices of degree $m$ and $41m$ vertices of degree 3.''
\end{quote}

\noindent
Our target is the \emph{first sentence} of the third sub-bullet under the bullet ``Related to
semicubic graphs'' in Section~6, ``Open questions'', of \cite{AAB}, on p.~29, verbatim
(subscripts written $H_{12}$, $T_{12}$, $H'_{12}$):

\begin{quote}
``Find different voltage assignments for the graph $H_{12}$ that produce graphs of girth 12 for
$m \in \{3, 4, \dots , 9\}$, or show no assignments exist. \emph{Preliminary investigations
suggest that new arc assignments are required to find smaller examples.} Does there exist a
different graph $H'_{12}$ (for example, using the same modified tree structure and arcs
$a, b, c, d$ but different arcs among the remaining leaves of $T_{12}$) that produces graphs of
girth 12 for these smaller values of $m$? For example, can we find graphs of small girth if we
construct a 16-cycle among the remaining leaves of $T_{12}$, or if we interlace two 8-cycles
differently?''
\end{quote}

The parent bullet, pp.~28--29, reads ``Related to semicubic graphs: Find voltage assignments to
construct graphs of girth 10 and 12 for missing values of $m$. More specifically:''. The
locator for the target is therefore: \cite{AAB}, Section~6 ``Open questions'', third
sub-bullet under ``Related to semicubic graphs'', p.~29.

\medskip
\noindent\textbf{Scope.}
That sub-bullet contains \emph{two} questions. Everything below concerns the first, and only in
the form it takes for the explicit graph of Section~\ref{sec:base}; nothing here bears on the
second, the different graph $H'_{12}$. Nothing here may be read as ``the third open question is
closed'': that reading needs the identification of Section~\ref{sec:corr}, which is not
established here.

\medskip
\noindent
Why $m=9$ was singled out by the authors, \cite[p.~28]{AAB}, verbatim:

\begin{quote}
``To compare the parameters we recall that the lower bound given in Lemma 3.4 in [6] is
$n(\{3,m\};12) \ge \lceil 109m/3\rceil + 17$. The graphs $(G_{12}, m)$ with $m \ge 9$ has order
$49m + 3$ whereas the graphs $(H_{12}, m)$ with $m \ge 10$ has order $41m + 3$. \dots
Interestingly, if we are able to find a voltage assignment to construct an $(H_{12},9)$ we
would be able to construct a graph of order 372.''
\end{quote}

\noindent
That order-$372$ graph is exhibited in Section~\ref{sec:exist}. The lower bound quoted there is
attributed by \cite{AAB} to Lemma~3.4 of Araujo-Pardo, Balbuena, Garc\'ia-V\'azquez, Marcote
and Valenzuela \cite{ABGMV}.

\section{The $44$-vertex base graph, in full}\label{sec:base}

Everything below is a statement about one explicit $44$-vertex graph, so that graph is printed
here completely: a reader needs neither the source's figure nor any program. It is our reading
of the graph of \cite[Figure~16]{AAB}; that reading, and the two places where the source's
printed tables have to be read against their own deletion list to obtain it, are recorded in
Section~\ref{sec:corr}, which also states what has \emph{not} been verified about the
identification. \emph{Convention.} From here on ``$H_{12}$'' abbreviates the graph printed in
this section, and nothing else; it is not asserted to be the graph of \cite[Figure~16]{AAB}.

\medskip
\noindent\textbf{The $44$ vertices} ($x^{*},y^{*},z^{*}$ are the three pinned vertices):
{\small
\begin{verbatim}
x*, x, x0, x1, x00, x01, x10, x11, x010, x011, x100, x101, x110,
x111, x0100, x0101, x0110, x0111, x1010, x1011, x1110, x1111,
y*, y, y0, y1, y01, y10, y11, y010, y011, y101, y110,
z*, z, z0, z1, z01, z10, z11, z010, z011, z101, z110
\end{verbatim}}

\noindent\textbf{The $43$ tree edges, all at voltage $0$:}
{\footnotesize
\begin{verbatim}
x*-x, x-x0, x-x1, x0-x00, x0-x01, x1-x10, x1-x11, x01-x010, x01-x011,
x10-x100, x10-x101, x11-x110, x11-x111, x010-x0100, x010-x0101,
x011-x0110, x011-x0111, x101-x1010, x101-x1011, x111-x1110, x111-x1111,
y*-y, y-y0, y-y1, y0-y01, y1-y10, y1-y11, y01-y010, y01-y011,
y10-y101, y11-y110,
z*-z, z-z0, z-z1, z0-z01, z1-z10, z1-z11, z01-z010, z01-z011,
z10-z101, z11-z110,
x00-y0, x00-z0
\end{verbatim}}

\noindent\textbf{The $20$ labelled arcs} (name: tail $\to$ head). These are the only edges
carrying a voltage, and they are the $20$ coordinates of every voltage tuple in this note,
always in this order:
{\footnotesize
\begin{verbatim}
a: x100 ->y10     b: x100 ->z10     c: x110 ->y11     d: x110 ->z11
e: x0100->y010    j: x0101->y011    q: x0110->y101    u: x0111->y110
w: x1010->y110    s: x1011->y101    l: x1110->y011    g: x1111->y010
f: x0100->z101    k: x0101->z110    r: x0110->z010    v: x0111->z011
al:x1010->z010    t: x1011->z011    p: x1110->z101    h: x1111->z110
\end{verbatim}}
\noindent(the arc written \texttt{al} is the source's $\alpha$). The coordinate order is
\[
 \varphi=[\,a,b,c,d,e,j,q,u,w,s,l,g,f,k,r,v,\alpha,t,p,h\,].
\]

\medskip
\noindent\textbf{The lift.} For $\varphi\colon\{\text{arcs}\}\to\mathbb{Z}_m$, the derived graph
$(H_{12},\varphi)$ has, for each non-pinned base vertex $v$, the $m$ vertices $(v,i)$,
$i\in\mathbb{Z}_m$, together with the three single vertices $x^{*},y^{*},z^{*}$. A tree edge
$u$--$v$ contributes $(u,i)$--$(v,i)$ for every $i$; an arc $X\colon t\to h$ contributes
$(t,i)$--$(h,\,i+\varphi(X)\bmod m)$ for every $i$; and $x^{*}$ is joined to $(x,i)$ for all
$i$, likewise $y^{*}$ to $(y,i)$ and $z^{*}$ to $(z,i)$.

\begin{proposition}\label{prop:base}
The graph above has $44$ vertices and $63$ edges, no loop and no repeated edge; its degree
sequence is $1$ three times and $3$ forty-one times, and the three vertices of degree~$1$ are
exactly $x^{*},y^{*},z^{*}$; the $43$ tree edges form a spanning tree, so the cycle rank is
$63-44+1=20$ and the $20$ arcs are a cotree; and it is bipartite, with classes of sizes $23$
and $21$, every one of the $20$ arcs joining opposite classes. Consequently $(H_{12},\varphi)$
has order $41m+3$ and $63m$ edges, exactly three vertices of degree $m$ and $41m$ of
degree~$3$, and is bipartite.
\end{proposition}

Restoring the eight leaves that \cite[p.~27]{AAB} deletes from $T_{12}$ to form $H_{12}$ ---
namely $x1000$, $x1001$, $x1100$, $x1101$, $y100$, $y111$, $z100$, $z111$, attached to
$x100,x100,x110,x110,y10,y11,z10,z11$ respectively --- gives a tree on $52$ vertices with $51$
edges, hence $49$ non-pinned vertices. That reproduces the source's own $49m+3$ and its
``$441=49\cdot 9$ vertices of degree 3'' for the companion family $G_{12}$, and $52-8=44$,
$51-8=43$ reproduce Theorem~5.6's ``$41m$ vertices of degree 3''.

\section{The search space, and the necessary condition}\label{sec:space}

\subsection*{Why $\mathbb{Z}_m^{20}$ with the tree at $0$ is the whole space}
Two independent reasons. First, it is literally the space the question asks about: in
\cite{AAB} the graph $H_{12}$ carries voltages on exactly the $20$ labelled arcs of Table~4,
and every other edge is unlabelled, i.e.\ at voltage~$0$. Second, it is gauge-complete anyway.
\emph{Switching} --- re-indexing the fibre over each base vertex by an arbitrary
$f\colon V\to\mathbb{Z}_m$ --- zeroes the voltages on any spanning tree without changing the
derived graph up to isomorphism. By Proposition~\ref{prop:base} the $43$ tree edges \emph{are}
a spanning tree and the cycle rank is $20$, so the normalised space has dimension exactly the
cycle rank. The pinned vertices do not break this: a pinned vertex is joined to \emph{all} $m$
copies of its unique neighbour, so the voltage on its pendant edge is irrelevant, and switching
at that neighbour merely permutes the fibre cyclically, which the pinned join is invariant
under. A census of $\mathbb{Z}_m^{20}$ therefore meets every isomorphism class of lift and
decides the cell as posed.

\subsection*{The short-cycle system}
A base cycle of length $q$ whose arc voltages sum to $\sigma$ lifts to cycles of length
$q\cdot m/\gcd(\sigma,m)$. Hence

\begin{lemma}\label{lem:nec}
If $(H_{12},\varphi)$ has girth $\ge 12$ then $\sigma_C(\varphi)\not\equiv 0 \pmod m$ for every
cycle $C$ of $H_{12}$ of length $<12$, where $\sigma_C$ is the sum of the arc voltages picked
up by traversing $C$ once, each arc counted $+1$ if traversed tail to head and $-1$ otherwise.
\end{lemma}

The condition is \emph{necessary only}, which is the safe direction for a nonexistence proof and
the only direction used in Section~\ref{sec:none}: unsatisfiability of the necessary system
rules out girth $\ge 12$, hence a fortiori girth exactly $12$. (Identifying the three pinned
fibres can only merge vertices, which never lengthens a cycle; closed walks and lollipops can
only lower the girth further.)

The census of cycles of $H_{12}$ shorter than $12$ is complete and finite: there are
\[
 18 + 27 + 82 = 127 \text{ cycles of lengths } 6,\,8,\,10,
 \quad\text{i.e. } \mathbf{254} \text{ directed},
\]
and there is no cycle of length below $6$ and none of odd length (the graph is bipartite). The
$254$ is the source's own published integer for $H_{12}$, \cite[p.~27]{AAB}: ``There are 254
such cycles''. All $127$ coefficient vectors have entries in $\{-1,0,+1\}$, with arity
histogram $\{1{:}16,\,2{:}44,\,3{:}4,\,4{:}55,\,6{:}6,\,8{:}2\}$. The $16$ arity-one
constraints force $16$ of the $20$ arcs to be nonzero; the four arcs \emph{not} so forced are
exactly $w,s,p,h$, each of whose single-arc fundamental cycle has length $12$.

\section{The existence half: three witnesses}\label{sec:exist}

\begin{theorem}\label{thm:exist}
For each $m\in\{7,8,9\}$ the voltage assignment $\varphi$ listed below makes $(H_{12},\varphi)$
a connected bipartite $(\{3,m\};12)$-graph on $41m+3$ vertices, with $63m$ edges, no loop and
no multiple edge, exactly three vertices of degree~$m$ and $41m$ of degree~$3$, and girth
exactly $12$.
\end{theorem}

\noindent In the coordinate order $[a,b,c,d,e,j,q,u,w,s,l,g,f,k,r,v,\alpha,t,p,h]$:

\medskip
\noindent$W_7$, a $(\{3,7\};12)$-graph on $\mathbf{290}$ vertices:
{\small
\begin{verbatim}
phi = [1, 2, 2, 1, 2, 4, 6, 3, 0, 5, 5, 6, 1, 5, 3, 1, 1, 4, 0, 3]
\end{verbatim}}
\noindent A shortest cycle, as lift vertices $(\text{base vertex},\text{fibre index})$:
{\small
\begin{verbatim}
(y11,5) (y1,5) (y,5) (y0,5) (y01,5) (y010,5)
(x0100,3) (z101,4) (z10,4) (z1,4) (z11,4) (x110,3)
\end{verbatim}}

\noindent$W_8$, a $(\{3,8\};12)$-graph on $\mathbf{331}$ vertices:
{\small
\begin{verbatim}
phi = [1, 2, 2, 1, 3, 7, 3, 1, 0, 4, 1, 7, 1, 3, 7, 3, 1, 7, 0, 4]
\end{verbatim}}
{\small
\begin{verbatim}
(x100,7) (y10,0) (y1,0) (y,0) (y0,0) (y01,0)
(y011,0) (x1110,7) (x111,7) (x11,7) (x1,7) (x10,7)
\end{verbatim}}

\noindent$W_9$, a $(\{3,9\};12)$-graph on $\mathbf{372}$ vertices --- the order the authors name
on their p.~28:
{\small
\begin{verbatim}
phi = [1, 2, 2, 1, 7, 6, 3, 1, 0, 5, 1, 3, 1, 3, 6, 7, 1, 3, 0, 5]
\end{verbatim}}
{\small
\begin{verbatim}
(x0,7) (x00,7) (z0,7) (z,7) (z1,7) (z10,7)
(z101,7) (x1110,7) (x111,7) (x11,7) (x1,7) (x,7)
\end{verbatim}}

\begin{proof}
Each witness is a finite exhibited object and every assertion of Theorem~\ref{thm:exist} is a
finite computation on it, carried out on the $41m+3$ vertex lift built from
Section~\ref{sec:base}: count the vertices and edges, check for loops and repeated pairs, take
the degree histogram, run one breadth-first search for connectivity, and run breadth-first
search from every vertex, capped at depth~$6$, for the girth. For a graph of girth $g\le 12$
there is a root on a shortest cycle from which the two arcs of that cycle meet at depth
$\lfloor g/2\rfloor\le 6$, so the cap loses nothing, and the minimum of
$d(u)+d(w)+1$ over the non-tree edges of all those searches is the girth exactly. The three
displayed $12$-cycles are the certificates that the girth is not larger: each is $12$ distinct
lift vertices whose $12$ consecutive pairs are edges of the lift, which one checks directly
against the tables of Section~\ref{sec:base}. That measurement is the only computation the
existence half depends on, and it uses no solver.
\end{proof}

\subsection*{Witnesses keeping the published values of $a,b,c,d$}
Theorem~\ref{thm:exist} is unrestricted. With $a,b,c,d$ fixed at the source's Table-4 values
$1,-1,-1,1$ and only the remaining $16$ arcs free, witnesses satisfying every assertion of
Theorem~\ref{thm:exist} still exist at all three values of $m$:
{\footnotesize
\begin{verbatim}
m=7  phi = [1, 6, 6, 1, 6, 5, 2, 5, 0, 5, 3, 2, 1, 6, 4, 6, 5, 4, 2, 5]
m=8  phi = [1, 7, 7, 1, 6, 5, 5, 1, 0, 3, 6, 1, 1, 4, 7, 5, 1, 2, 0, 2]
m=9  phi = [1, 8, 8, 1, 1, 6, 3, 1, 0, 7, 1, 6, 2, 4, 5, 8, 1, 2, 0, 3]
\end{verbatim}}
\noindent (here $6=-1$ in $\mathbb{Z}_7$, $7=-1$ in $\mathbb{Z}_8$, $8=-1$ in $\mathbb{Z}_9$).
For $m\le 6$ the restricted cell is empty a fortiori, since by
Theorem~\ref{thm:none} the full cell is.

\section{The nonexistence half: $m=3,4,5,6$}\label{sec:none}

\begin{theorem}\label{thm:none}
For each $m\in\{3,4,5,6\}$ there is \emph{no} $\varphi\in\mathbb{Z}_m^{20}$ for which
$(H_{12},\varphi)$ has girth $\ge 12$; in particular none for which it has girth exactly $12$.
\end{theorem}

\begin{proof}
By Lemma~\ref{lem:nec} it suffices to exhibit, for each $m$, a set of arcs and a set of
short-cycle constraints supported on those arcs alone which is already unsatisfiable over
$\mathbb{Z}_m$: a partial assignment violating a necessary constraint cannot be completed, so
the whole cell $\mathbb{Z}_m^{20}$ is decided. The following four sets work; in each case
\emph{every} tuple over the listed arcs is enumerated flat, with no search and no propagation,
and none survives.
\[
\begin{array}{c|l|r|r}
 m & \text{arcs} & \text{tuples} & \text{survivors}\\\hline
 3 & g,f,k,p,h & 3^{5}=243 & 0\\
 4 & u,w,s,r,v,\alpha,t & 4^{7}=16{,}384 & 0\\
 5 & q,u,w,s,r,v,\alpha,t & 5^{8}=390{,}625 & 0\\
 6 & q,u,w,s,r,v,\alpha,t & 6^{8}=1{,}679{,}616 & 0
\end{array}
\]
The constraints used are, in each case, exactly those of the $127$ of Section~\ref{sec:space}
whose support lies inside the listed arc set; they are recomputed from the tables of
Section~\ref{sec:base} by the accompanying program. No enumeration over $\mathbb{Z}_m^{20}$
itself is performed or needed: what is enumerated is only the four small cores in the table, and
the orders $41m+3$ that the four empty cells would have given are $126$, $167$, $208$ and $249$.
\end{proof}

\subsection*{$m=3$ is a pencil-and-paper proof}
For $m=3$ the ten constraints supported on $g,f,k,p,h$ read, over $\mathbb{Z}_3$,
\begin{gather*}
 g\ne 0,\quad f\ne 0,\quad k\ne 0,\quad f\ne k,\quad p\ne f,\quad p\ne h,\\
 h\ne g,\quad h\ne k,\quad h\ne g+k,\quad f-k-p+h\ne 0 .
\end{gather*}
(For instance $f-k-p+h\ne0$ comes from the octagon
$x0100 \xrightarrow{f} z101 \xleftarrow{p} x1110 - x111 - x1111 \xrightarrow{h} z110
\xleftarrow{k} x0101 - x010 - x0100$.) Since $g,f,k$ are nonzero they lie in $\{1,2\}$, and
$f\ne k$ forces $\{f,k\}=\{1,2\}$. Four cases:

\begin{itemize}
\item $g=1,k=1$ (so $f=2$): $h\ne 1$ and $h\ne g+k=2$, so $h=0$; then $p\ne f=2$ and
      $p\ne h=0$, so $p=1$; but $f-k-p+h=2-1-1+0=0$. Contradiction.
\item $g=1,k=2$ (so $f=1$): $h\ne 1$, $h\ne 2$, $h\ne g+k=0$ --- no value remains.
\item $g=2,k=1$ (so $f=2$): $h\ne 2$, $h\ne 1$, $h\ne g+k=0$ --- no value remains.
\item $g=2,k=2$ (so $f=1$): $h\ne 2$ and $h\ne g+k=1$, so $h=0$; then $p\ne 1$ and $p\ne 0$, so
      $p=2$; but $f-k-p+h=1-2-2+0=-3\equiv 0$. Contradiction.
\end{itemize}

\noindent
So the $m=3$ half needs no computer. At $m=3$ the three pinned vertices have degree $m=3$ as
well, so the lift is \emph{cubic}, on $41\cdot3+3=126$ vertices, and $126$ is the order of the
$(3,12)$-cage. A cubic graph of girth $12$ on $126$ vertices \emph{is} a $(3,12)$-cage, and the
$(3,12)$-cage is unique \cite{DS16}. Hence

\begin{corollary}\label{cor:cage}
The $(3,12)$-cage is not a $\mathbb{Z}_3$ lift of $H_{12}$; equivalently, $H_{12}$ does not
present the Tutte $12$-cage \textup{\cite{Tutte}} as a $\mathbb{Z}_3$ lift.
\end{corollary}

\begin{remark}
Two riders on Corollary~\ref{cor:cage}. (i) The \emph{uniqueness} of the $(3,12)$-cage is
quoted from the literature \cite{DS16} and was not re-verified here; what is proved here is the
emptiness of $\mathbb{Z}_3^{20}$, plus the arithmetic $41\cdot 3+3=126$. (ii) The notation
$n(\{3,3\};12)$ is not used anywhere in this note and should not be: the framework in which
every other number here is stated requires the two degrees to be distinct, so the $m=3$ cell
lies outside it. Write ``the $(3,12)$-cage'', or ``$(H_{12},3)$''.
\end{remark}

\begin{remark}[the boundary of Theorem~\ref{thm:none}]
Theorem~\ref{thm:none} is a statement about $20$-arc voltage assignments on the graph of
Section~\ref{sec:base}, in the switching normal form of Section~\ref{sec:space}. It does
\emph{not} say $n(\{3,m\};12)>41m+3$ for those $m$, and nothing here may be read that way. It is
a nonexistence statement about one construction.
\end{remark}

\section{Consequences for the published bounds}\label{sec:record}

The published record table for bi-regular cages is compiled by Goedgebeur, Jooken and
Van~den~Eede \cite{GJV}; the four rows relevant here read, in their version of record,
\[
\begin{array}{c|r|r}
 m & \text{lower bound} & \text{upper bound}\\\hline
 6 & 235 & 250 \to 243\\
 7 & 272 & 334\\
 8 & 308 & 374\\
 9 & 344 & 374 \to 360
\end{array}
\]
where ``$x\to y$'' is their own notation for ``$x$ from the literature, $y$ our improvement''.
Theorem~\ref{thm:exist} therefore gives

\begin{corollary}\label{cor:record}
$n(\{3,7\};12)\le 290$ and $n(\{3,8\};12)\le 331$.
\end{corollary}

Both lie strictly inside the published feasible window --- $\lceil 109\cdot7/3\rceil+17=272<290$
and $\lceil 109\cdot 8/3\rceil+17=308<331$ --- so the remaining gaps are $18$ and $23$.

\medskip
\noindent\textbf{Attribution.} The $334$ at $m=7$ belongs to Araujo-Pardo, Balbuena,
L\'opez-Ch\'avez and Montejano \cite{ABLM}, Theorem~2.2 --- a reference of \cite{AAB} itself,
not a bound of \cite{AAB}: the source's own Table~1 (p.~4) prints
$\mathrm{rec}(\{3,7\};12)=376$, i.e.\ it does not cite the better $334$ that sits in its own
bibliography. So no sentence here improves the source's authors at $m=7$. The $374$ at $m=8$
\emph{is} theirs (Corollary~2.4 of \cite{AAB}); only there does this note improve a bound of
\cite{AAB}.

\medskip
\noindent\textbf{Adjacent work, and what it does not settle.} \cite{GJV} is the current
compilation of the $n(\{r,m\};g)$ bounds and the nearest work to this note: its Section~4
generalises the \emph{gluing} theorem of \cite{AAB} --- ``In [2] Aguilar, Araujo-Pardo, and
Berman generalized Theorem 3 from [5]\dots\ we generalize their theorem a) for all girths
$g\ge5$, b) for more pairs $(r,m)$, c) by improving their bound'' --- and not the $H_{12}$
voltage-graph construction. It therefore says nothing about voltage assignments on $H_{12}$ for
$m<10$, and its girth-$12$ rows leave the $m=7$ and $m=8$ upper bounds $334$ and $374$ with an
empty improvement column; its own improvements in this range are at $m=6$ ($250\to243$) and
$m=9$ ($374\to360$), the latter being the $360$ that keeps the $m=9$ witness below from carrying
a record claim. \cite{ABLM}, Theorem~2.2, supplies the $334$ by a construction unrelated to
$H_{12}$ and settles nothing about it; \cite{ABGMV}, Lemma~3.4, supplies only the lower bound
$\lceil 109m/3\rceil+17$ and no upper bound. None of the three addresses the sub-bullet's
question, and none of them settles the identification of Section~\ref{sec:corr}.

\medskip
\noindent\textbf{The $m=8$ bound against Table~1 of the source.}
Table~1 of \cite{AAB}, p.~4, prints $\mathrm{rec}(\{3,8\};12)=\mathbf{304}$, which is
\emph{below} our $331$, so a referee reading that page will read
Corollary~\ref{cor:record} as false; the claim should not be quoted without this paragraph. We
read that $304$ as a misprint for $374$. No erratum is available to us and none is verified
here; what is checked is the following. (i) $304<308$, and $308$ is the lower bound printed in
the adjacent column of the same row, reproduced by
$\lceil 109\cdot8/3\rceil+17=291+17=308$: so the printed $304$ is inconsistent with the lower
bound of its own row. (ii) Table~1's technique column names Corollary~2.4 of \cite{AAB} (p.~7)
for that entry: ``There exist $(3,m;12)$-graphs of order: $41m+\frac{m}{3}+2$ for $m=3k$;
$41m+\frac{m}{3}+86+\frac{2}{3}$ for $m=3k+1$; $41m+\frac{m}{3}+43+\frac{1}{3}$ for $m=3k+2$'',
and at $m=8=3\cdot2+2$ that is $328+\frac83+43+\frac13=374$, while in exact rational arithmetic
the same corollary reproduces every other $\mathrm{rec}(\{3,m\};12)$ entry of Table~1 ($252$,
$250$, $250$, $376$, $374$ at $m=4,5,6,7,9$). (iii) \cite{GJV} print $374$ for that row,
attributed to the same corollary, with an empty improvement column.

\medskip
\noindent\textbf{The $m=9$ witness carries no record claim.}
$\;372>360$, and $360$ is published \cite{GJV}. The $372$ of Theorem~\ref{thm:exist} realises
the order the authors of \cite{AAB} name on their p.~28, but it improves nothing, and no
sentence here may say that it does. The orders $413$ at $m=10$ and $454$ at $m=11$ are
Theorem~5.6 of \cite{AAB} and not ours. And $m=6$ carries nothing either:
\cite{GJV}'s improved $n(\{3,6\};12)\le 243$ is below the order $249$ that $(H_{12},6)$ would
have had --- moot in any case, since Theorem~\ref{thm:none} empties that cell.

\section{Reading $H_{12}$ off the source's tables}\label{sec:corr}

Tables~3 and~4 of \cite{AAB}, and the prose on its p.~27, cannot be transcribed literally into a
graph whose $41$ non-pinned vertices all have degree~$3$: they name leaves that the same page
deletes. Section~\ref{sec:base} therefore prints the graph we read from
\cite[Figure~16]{AAB}, and the two places where that reading departs from the printed tables are
these.

\begin{enumerate}
\item \textbf{Two endpoint misprints, against the source's own deletion list.} Tables~3 and 4
print the $x$-leaf sequence ``\dots $x1010$, $x1001$ \dots'' where the surviving leaf set forces
$x1011$ ($x1001$ is itself one of the eight deleted leaves, and is used twice while $x1011$ is
never used); and Table~4's third row gives starting leaves $x1100,x1101$ for the arcs $p,h$,
although both are deleted on p.~27, the surviving leaves still needing a $z$-arc being
$x1110,x1111$.
\item \textbf{The p.~27 prose ``we added new labeled directed arcs from $x100$ to $y10$, $z10$
and from $x110$ to $y111$ and $z111$'' cannot be read literally}, since $y111$ and $z111$ are
both in the deletion list; we read it as Table~4's $x110\to y11,z11$, which is what gives
$y11,z11$ degree~$3$.
\end{enumerate}

\noindent\textbf{What supports the identification, and what does not.} On the reading above all
$41$ non-pinned vertices are cubic, and the short-cycle census of Section~\ref{sec:space}
reproduces the source's own published $254$ and its $41m+3$; restoring the eight deleted leaves
reproduces its $49m+3$ and $441$. Those are the whole of the evidence offered here. We do
\emph{not} supply an independent, arc-by-arc transcription of Figure~16, and the accompanying
program does not check one: it is given the graph of Section~\ref{sec:base} and checks
statements about that graph. Nor do we claim that the reading above is the only reading with all
$41$ vertices cubic. Accordingly Proposition~\ref{prop:base},
Theorems~\ref{thm:exist} and~\ref{thm:none} and Corollaries~\ref{cor:cage}
and~\ref{cor:record} are statements about the graph printed in Section~\ref{sec:base}, and this
note does not claim more. In particular it does not claim to answer the question of \cite{AAB}:
that would need the printed graph to be the $H_{12}$ of \cite[Figure~16]{AAB}, which is exactly
what is not established here and what a referee would have to check against the source, arc by
arc. Corollary~\ref{cor:record} is the one consequence unaffected by this, since its two
witnesses are exhibited and checked as graphs in their own right.

\section{What is not settled}\label{sec:open}

\begin{itemize}
\item \textbf{The identification of the printed graph with \cite[Figure~16]{AAB} is not
established} (Section~\ref{sec:corr}), so the question of \cite{AAB} is not claimed to be
answered, in either direction.
\item \textbf{The second question of the sub-bullet --- the different graph $H'_{12}$ --- is
completely untouched and remains open, at every $m$.} Nothing here rewires the leaf-joining arcs
of $T_{12}$, builds a $16$-cycle among the remaining leaves, or interlaces two $8$-cycles
differently, and nothing here says whether such an $H'_{12}$ exists at any $m$.
\item \textbf{The \emph{second} sub-bullet of the same page is a different question and is not
answered here.} Verbatim, \cite[p.~29]{AAB}: ``Find voltage assignments to produce
$(G_{12};m)$ graphs with the construction given in Theorem 5.5 for any $4\le m\le 8$, or show
no such graphs exist''. That is the $G_{12}$ family of order $49m+3$, not the $41m+3$ family
treated here.
\item \textbf{Optimality is open.} Whether $290$ and $331$ are the true values of
$n(\{3,7\};12)$ and $n(\{3,8\};12)$ is not addressed; the published lower bounds $272$ and
$308$ stand untouched. No attempt was made to raise them, and none to find a smaller
$(\{3,7\};12)$ or $(\{3,8\};12)$ graph outside this family.
\item \textbf{Counting is not claimed}, at any $m$: the number of assignments passing the
necessary system of Lemma~\ref{lem:nec} is not the number giving girth $12$.
\item \textbf{Isomorphism questions are unexamined:} whether the witnesses of
Section~\ref{sec:exist} are pairwise non-isomorphic, and whether any is isomorphic to a known
graph, was not checked.
\item \textbf{The uniqueness of the $(3,12)$-cage is an inherited external fact}, used only in
Corollary~\ref{cor:cage} and flagged there.
\item \textbf{Nonexistence is about this construction, not about the cage numbers} --- see the
remark closing Section~\ref{sec:none}. In particular the emptiness at $m=4,5,6$ says nothing
about \cite{GJV}'s $n(\{3,4\};12)\le 220$, $n(\{3,5\};12)\le 230$, $n(\{3,6\};12)\le 243$.
\end{itemize}

\section{Verification}\label{sec:verify}

Every computational assertion of Sections~\ref{sec:base}--\ref{sec:record} is re-derived, from
the objects printed above and from nothing else, by the accompanying program \texttt{verify.py}
(Python~3.9+, standard library
only, exact integer and \texttt{Fraction} arithmetic, no floating point in any decision). It
rebuilds the base graph from the tables of Section~\ref{sec:base} and checks
Proposition~\ref{prop:base}; recomputes the short-cycle census and the arity histogram of
Section~\ref{sec:space}, reproducing the source's published $254$; builds all six lifts and
measures order, edge count, loops, repeated pairs, degree histogram, connectivity and exact
girth, and checks the three displayed $12$-cycles edge by edge; re-runs the four flat
enumerations of Theorem~\ref{thm:none} and the four-case argument for $m=3$; and re-derives the
arithmetic of Section~\ref{sec:record}, including Corollary~2.4 of \cite{AAB} in exact
rationals at $m=4,\dots,9$. Its recorded run is in
\texttt{verify.output.txt}. Its output
also carries arithmetic side points this note does not state --- among them the total size of the
four cells emptied by Theorem~\ref{thm:none}, and evaluations of \cite{GJV}'s Theorem~4.2 and of
Table~1 of \cite{AAB}; the Theorem~4.2 checks evaluate selected parameter choices and are not
minimisations over its cases and parameters, and no check of the run verifies an erratum. The
program
is a convenience, not a dependency: the existence half reduces to one breadth-first search on a
$290$-vertex graph, and the $m=3$ half is the four-case argument above.

\medskip
\noindent\textbf{What \texttt{verify.py} does not cover.} It does not check that the graph of
Section~\ref{sec:base} is the graph of \cite[Figure~16]{AAB}: the program is handed that graph,
and the readings of Section~\ref{sec:corr} are assertions about the source's tables, which are
not reprinted here. A reader who wants them checked must fetch Tables~3 and~4 and Figure~16 from
\cite{AAB}. Beyond that, the program's own closing scope note lists the inherited external facts
it does not re-run, among them the uniqueness of the $(3,12)$-cage used in
Corollary~\ref{cor:cage} and the record values $334$, $374$ and $360$, which are transcribed
from the literature and not recomputed.

\begin{thebibliography}{9}

\bibitem{AAB}
F.~Aguilar, G.~Araujo-Pardo, and L.~W. Berman,
\emph{Semicubic cages and small graphs of even girth derived from voltage graphs},
Ars Math. Contemp. \textbf{26} (2026), \#P3.03.
\href{https://doi.org/10.26493/1855-3974.3116.35e}{doi:10.26493/1855-3974.3116.35e};
preprint \href{https://arxiv.org/abs/2305.03290}{arXiv:2305.03290}.
All page references are to the $30$-page journal galley.

\bibitem{GJV}
J.~Goedgebeur, J.~Jooken, and Van~den~Eede,
\emph{Computational methods for finding bi-regular cages},
Ars Math. Contemp., published online 15 December 2025.
\href{https://doi.org/10.26493/1855-3974.3497.56c}{doi:10.26493/1855-3974.3497.56c};
preprint \href{https://arxiv.org/abs/2411.17351}{arXiv:2411.17351}.

\bibitem{ABLM}
Araujo-Pardo, Balbuena, L\'opez-Ch\'avez, and Montejano,
\emph{On bi-regular cages of even girth at least 8},
Aequationes Math. \textbf{86} (2013), 201--216.

\bibitem{ABGMV}
Araujo-Pardo, Balbuena, Garc\'ia-V\'azquez, Marcote, and Valenzuela,
Discrete Math. \textbf{308} (2008), 2484--2491.
The source of the lower bound $n(\{3,m\};12)\ge\lceil 109m/3\rceil+17$, as attributed by
\cite[Lemma~3.4 of its reference {[6]}, quoted p.~28]{AAB}.

\bibitem{Tutte}
W.~T. Tutte,
\emph{A family of cubical graphs},
Proc. Cambridge Philos. Soc. \textbf{43} (1947), 459--474.

\bibitem{DS16}
G.~Exoo and R.~Jajcay,
\emph{Dynamic cage survey},
Electron. J. Combin., Dynamic Survey DS16.
Quoted here only for the order and the uniqueness of the $(3,12)$-cage.

\end{thebibliography}

\end{document}
